\section{The Depth Rubric}

The rubric is the central standard of the whole program. It exists to answer one question about every experiment we run. What did this test actually establish? Not whether it printed PASSED, not whether a green checkmark appeared, but what claim about the world it has earned the right to make.

The commitment behind the rubric is plain.

\begin{principle}[The honesty floor]
A wall of green checkmarks is worth less than three honest results with two failures. We grade our own tests as harshly as a stranger's, because the easiest person to fool is ourselves.
\end{principle}

A test that clears its threshold has earned a \texttt{pass}. It has not earned the label \emph{discovery}. The two are different kinds of fact, and the depth grade is what tells them apart. The grade is stated with the result, every time, so a reader never has to guess what a number means.

The rubric works alongside the experimental architecture of the previous section. That section gives the runner that enforces these grades. This section gives the grades themselves, the second axis of status tags that runs beside them, and the disciplines that keep an imposed answer from being read as an emergent one.

\subsection{The four depth levels}

Every experiment is graded on one of four levels. The level names the \emph{kind} of result, ordered from the worst failure of method up to the genuine target.

\begin{table}[ht]
\centering
\begin{tabular}{@{}llll@{}}
\toprule
level & name & what it establishes & example \\
\midrule
L0 & circular & the answer is put in by hand & integrate a hardcoded equation, report the constant inside it \\
L1 & known math & confirms an established mathematical fact & the 24 sites form a flat 4D lattice \\
L2 & known physics & reproduces a known construction here & the Dirac quantum walk on the mesh \\
L3 & emergent and novel & the knit produces it as a measured consequence & the genuine target \\
\bottomrule
\end{tabular}
\caption{The four depth levels. The grade is declared with every result and never inferred by a reader.}
\label{tab:depth-levels}
\end{table}

\subsubsection{L0, circular}

The conclusion is the input restated. This is the worst and the most common failure of method, and the rubric names it first so it is never confused with a result.

\begin{definition}[A circular check]
A check is \emph{circular} when the quantity it reports was set by hand somewhere in its own machinery. The test fixes a constant, runs an equation that contains that constant, and then reads the constant back out of the output.
\end{definition}

The shape is easy to see once stated. Set a coupling to $2$, integrate the equation that already holds the $2$, then announce that the coupling is $2$. Nothing about the dynamics has been tested. The output was guaranteed by the input. It is a tautology dressed as proof.

L0 results are not deleted in silence. They are kept, labeled circular, and quarantined. The catalog holds fifteen of them, sorted to the bottom, never cited as evidence. A check that turns out circular under audit is relabeled as a consistency note, not counted as a finding. The relabel is the honest move. A consistency note still has value, it confirms that two parts of the construction agree, but it is not evidence that the knit produced anything.

\subsubsection{L1, known math}

The test correctly confirms an established mathematical fact about the structure. These results are real and correct. They are not discoveries.

\begin{result}[The math floor]
An L1 result establishes a true and previously known mathematical fact about the mesh, the docks, or the sites. It is labeled L1 so that no one mistakes a property known in advance for an emergent one.
\end{result}

The examples on this substrate are clean. The 24 sites of a dock are the $D_4$ root system, which is a flat four-dimensional lattice. The dock geometry is the 24-cell, which is the binary tetrahedral group. A $2\pi$ rotation on the sites carries a minus sign. Each of these is true, and each was true before any experiment ran. The minus sign under a full turn is a property of the group, fixed by the algebra, not a thing the knit had to produce.

L1 is the math floor of the program. The \texttt{foundations} category and much of the substrate survey live here. They are correct, they are cited where they reproduce known work, and they are labeled L1 so the foundation is honest before anything is built on it. One hundred and fourteen experiments sit at this level.

\subsubsection{L2, known physics}

The test reproduces a known physical construction on this substrate. The construction is already established in the literature. The contribution is showing that it runs on the collide-and-stream knit and the $\{3,4,3,4\}$ mesh, not that the construction itself is new.

\begin{result}[Reproduction, honestly labeled]
An L2 result reproduces a known physical construction on the mesh under the knit. The work is real and useful. The construction is cited to its source, and the result is labeled L2 so that reproducing known physics is never dressed as discovering new physics.
\end{result}

The examples are the Dirac quantum walk, lattice gauge theory carried on the mesh, and a ballistic light cone arising from the streaming step. All three are established. We cite them. Reproducing known work on a fresh substrate is honest and valuable, and it is exactly what a young program should spend most of its effort on.

L2 is the largest tier, two hundred and eighty-four of the four hundred and fifty-seven experiments. A young program reproducing known physics on a new substrate is doing useful work. The label is what keeps it honest.

\subsubsection{L3, emergent and novel}

This is the genuine target. The knit produces the result as a measured consequence, with a control, and ideally with a quantitative prediction that could have come out wrong.

The bar has three parts, and all three must be met.

\begin{enumerate}
\item \textbf{Measured, not assumed.} A dispersion, a coupling, or a charge is read out of the actual dynamics on the mesh. It is not written in as a formula and then confirmed.
\item \textbf{A control that could have failed.} There is a case where the answer should be NO, run through the same machinery, and the test gives NO there.
\item \textbf{Quantitative where possible.} A number that could have come out wrong is stronger than a boolean yes.
\end{enumerate}

The contrast with L1 is sharp and worth stating outright. The algebra fact that a $2\pi$ rotation on the sites carries a minus sign is L1, a property of the group known in advance. The dynamical version, a spinor that carries that same sign as a measured consequence of the knit acting beat by beat, is the L3 target. The same minus sign appears in both, but only one of them was earned by the dynamics.

\begin{principle}[The rarity is the integrity]
L3 is small by design, forty-four of four hundred and fifty-seven experiments, about one in ten. A real emergent claim is rare and expensive. A program reporting mostly L3 results would be the suspicious one. It would mean the bar is too low or the grades are inflated.
\end{principle}

\subsection{The status tags}

The depth level grades the \emph{kind} of result. A second axis grades \emph{how} the result was established. The two axes are independent, and both are stated with every claim. They do not collapse into one.

Every claim carries a status tag.

\begin{table}[ht]
\centering
\begin{tabular}{@{}ll@{}}
\toprule
tag & meaning \\
\midrule
DERIVED & follows by argument or exact computation from the base atoms \\
MEASURED & computed on the substrate with a control \\
PREDICTED & a quantitative consequence implied but not yet checked \\
OPEN & genuinely unsettled \\
\bottomrule
\end{tabular}
\caption{The four status tags. Depth says how strong the kind of result is. The tag says how it was established.}
\label{tab:status-tags}
\end{table}

The two axes stay apart because they answer different questions. A DERIVED claim can be L1, an algebra fact about the sites that follows by exact computation, or it can sit under an L3 target as the argument half of an emergent claim. A MEASURED claim is the usual home of an L3 result, because L3 demands a measured consequence with a control. A PREDICTED claim names a number the framework implies that no run has yet returned. An OPEN claim names a question with no settled answer.

The scoreboard of the program reports DERIVED, MEASURED, PARTIAL, and OPEN. The word PARTIAL is the scoreboard's name for a claim whose control or quantitative half is not yet in place. That is the same gate that downgrades a controlless L3, described below. So the scoreboard tags and the depth levels are not two systems. They are one system read from two sides.

\subsection{The control requirement}

This is the rule that separates L3 from everything below it. It is the single hardest demand the rubric makes.

\begin{principle}[The control requirement]
An L3 claim must carry a control that could have failed. A control is a substrate, a knit, or a parameter where the answer should be NO, run through the same machinery, that gives NO.
\end{principle}

A test that cannot fail proves nothing. A yes with nothing that could have been a no is not evidence. It is the contrast between the positive case and the negative control that carries the meaning.

The canonical pair in this program is the spinor test.

\begin{table}[ht]
\centering
\begin{tabular}{@{}ll@{}}
\toprule
result & role \\
\midrule
the $\{5,3,4\}$ mesh carries no spinor & the control, the case that gives NO \\
the $\{3,4,3,4\}$ mesh carries a spinor & the claim, the case that gives YES \\
\bottomrule
\end{tabular}
\caption{The canonical control pair. The negative case is what makes the positive case mean something.}
\label{tab:spinor-control}
\end{table}

The reason this pair works is the decomposition of the directions. The 24 sites of a $\{3,4,3,4\}$ dock are the $D_4$ root system, which splits as $8v + 8s + 8c$, a vector plus two genuine half-integer-spin spinors. The 12 directions of $\{5,3,4\}$ split under icosahedral symmetry as $1 + 3 + 3' + 5$, all integer spin, with no spinor anywhere in the list. So the negative control is not a contrivance. It is a real substrate that genuinely cannot carry the thing the positive case carries.

\begin{result}[The control is the evidence]
Without the negative control, the $\{3,4,3,4\}$ spinor result is just a yes with nothing that could have been a no. With it, the contrast between the substrate that carries a spinor and the one that cannot is the evidence.
\end{result}

Controls live alongside the experiments as their own shared null cases, and the comparison substrates come from the tessellation catalog. An experiment imports the null case rather than rolling its own, so the same control is used everywhere it applies and never reinvented.

\subsection{The integrity rule}

The control requirement is enforced by the runner, not left to good intentions. The architecture holds this gate so that no author can talk past it.

\begin{principle}[The integrity rule]
A controlless L3 is downgraded to \texttt{partial}, and the build hard-fails on any \texttt{partial}. An experiment cannot declare \texttt{depth: 'L3'} and skip the control. The machinery catches it.
\end{principle}

The author has exactly two honest choices, and there is no third.

\begin{enumerate}
\item Supply the control, and the claim stands as L3.
\item Drop the control, and the result is L2, a reproduction, not a discovery.
\end{enumerate}

There is no path where a deep claim stands on a built-in answer. The runner that holds this gate is described in the section on the experimental architecture. The point here is that the gate is mechanical. It does not depend on the author remembering to be honest. The build will not pass until the control is present or the claim is downgraded.

\subsection{The honest-negative discipline}

Every experiment must be able to fail. A \texttt{fail} verdict is not a build break. It is a recorded scientific outcome, published like any other.

\begin{principle}[Every negative narrows the search]
A recorded \texttt{fail} is a finding. An imposed ingredient must never be read as emergence. If a phenomenon needs a further ingredient, that need is itself the finding.
\end{principle}

This discipline has a hard edge, and the edge is the base count. The model is exactly the eight atoms.

\begin{table}[ht]
\centering
\begin{tabular}{@{}ll@{}}
\toprule
atom & what it is \\
\midrule
mesh & the growing $\{3,4,3,4\}$ crystal as a whole \\
dock & one geometric cell, a here-now \\
site & one of the 24 directions out of a dock \\
tone & the ternary vibe value $\{-1, 0, +1\}$ \\
lean & the value order $-1 < 0 < +1$ \\
knit & the collide-then-stream law of the state \\
wake & the growing edge where new docks are born \\
beat & the discrete step \\
\bottomrule
\end{tabular}
\caption{The eight atoms. Nothing else is base. The arrow of time and the attraction are emergent, not base.}
\label{tab:eight-atoms}
\end{table}

That is the whole base. The arrow, the direction of time, is emergent, it comes from the wake acting through the lean, not from a clock written into the law. The attraction is emergent too. If a phenomenon needs cohesion, planning, a stabilizer, or a gravity term, that need is the finding. It goes into the notes as a negative. It is not a knob to turn quietly until the test goes green.

\begin{result}[The honest negative, stated]
If a phenomenon does not come from the eight atoms, we report the honest negative. We do not impose the missing ingredient and call the result emergent.
\end{result}

The largest example is the \texttt{selves} category. Many of its experiments are honest negatives, cases where a bound self did not emerge from the pure eight atoms alone. The pure knit churns and produces no self. That negative is published, not hidden, because the published negative is exactly where the integrity of the program lives.

\subsection{Why most results are L1 and L2}

A young program is mostly confirmation. That is the honest shape, and the catalog distribution shows it plainly.

\begin{table}[ht]
\centering
\begin{tabular}{@{}lrr@{}}
\toprule
depth & count & share \\
\midrule
L0 circular & 15 & about 3 percent \\
L1 known math & 114 & about 25 percent \\
L2 known physics & 284 & about 62 percent \\
L3 emergent and novel & 44 & about 10 percent \\
\bottomrule
\end{tabular}
\caption{The catalog distribution across the four hundred and fifty-seven experiments. The pyramid is broad at L1 and L2, narrow at L3.}
\label{tab:depth-distribution}
\end{table}

The reasoning is plain. Before a substrate can earn a novel claim, it has to earn the right to the next claim. It must first show that its math is correct, which is L1. It must then show that it reproduces the known physics, which is L2. Only on top of that foundation does an L3 claim mean anything.

\begin{principle}[The honest pyramid]
Most honest results in a young program are L1 and L2, and that is fine, as long as they are labeled as such. The pyramid, broad at the math and reproduction floors, narrow at the emergent peak, is what an honest measurement program looks like.
\end{principle}

\subsection{The grade is stated, never inferred}

The rubric only works if the grade is declared with the result, not guessed by a reader. So every experiment carries its depth in the registry, the catalog records it, and the printed claim is always read against it.

\begin{itemize}
\item PASSED at L1 means correct and established, not a discovery.
\item PASSED at L2 means a known construction reproduced here, not new.
\item PASSED at L3 means a measured emergent consequence, with a control that could have failed.
\end{itemize}

The discipline reduces to one refusal.

\begin{principle}[The one thing we never do]
The one thing we never do is fool ourselves. A result is worth reporting when it could have come out the other way and did not, when there is a control that could have failed and did not, and when an independent reader running the same command gets the same number. Until then it is a consistency check or a work in progress, and the rubric makes us say so.
\end{principle}

The companion section on the experimental architecture gives the experiment constructor, the verdict shape, the generated catalog, and the runner that enforces the control requirement. The standing methodology note holds the standard in full, with the anti-patterns and the double-and-triple-check rule that every test passes through before it is called a result.
